\section{Consolidated Status of the Three-Channel Program}
\label{sec:ConsolidatedStatus}

This section consolidates the results of Sections~\ref{sec:GreenTraceCalibration}--\ref{sec:CylinderRayleigh} into a single status theorem that states precisely what is proven, what is conditional, and what is open. The goal is honest bookkeeping: a reader should be able to determine from this section alone the exact logical state of the Jang-equation route to the spacetime Penrose inequality.

\subsection{Setup}

Let $(M,g,k)$ be a smooth asymptotically flat initial data set satisfying the dominant energy condition $\mu\ge|J|_g$. Let $\Sigma\subset M$ be a closed future-trapped surface ($\theta^+[\Sigma]\le0$). The Jang-equation reduction (Section~\ref{app:BK_Identity}) produces a sealed Jang manifold $(\widetilde M,\widetilde g)$ with cylindrical ends at the MOTS blow-up loci; assume $\Sigma$ is one such strictly stable MOTS blow-up with stability operator $L_\Sigma=-\Delta_\Sigma+|A|^2+\Ric(\nu,\nu)$ and $\lambda_1(L_\Sigma)>0$. Let $G_\infty$ be the ADM Green potential of Definition~\ref{def:ADMGreenNormalization37} with trace $\Gamma_\Sigma=G_\infty|_\Sigma$.

\subsection{Consolidated status theorem}

\begin{theorem}[Consolidated status]\label{thm:ConsolidatedStatus}
Under the setup above, the spacetime Penrose inequality $\sqrt{A(\Sigma)/16\pi}\le m_{\mathrm{ADM}}(g)$ holds provided any one of the following three conditions is satisfied:
\begin{enumerate}
    \item[\textup{(C1)}] \emph{Pointwise favorable jump (Section~\ref{app:LogicalStructure}):} $\tr_\Sigma k\ge0$ pointwise on $\Sigma$;
    \item[\textup{(C2)}] \emph{Drift-stability inequality (Section~\ref{sec:DriftAnchor}, Section~\ref{sec:VanishingSet}):} Hypothesis~\ref{hyp:DriftStability} holds, i.e.\ there exists $\delta_{\mathbf b}<1/2$ such that
    \[
        \biggl|\int_\Sigma\frac{\tr_\Sigma k}{|\beta_\Sigma|^2}\,\langle\beta_\Sigma,\nabla_\Sigma\Gamma_\Sigma^{\mathbf b}\rangle\,dA\biggr|
        \;\le\;\delta_{\mathbf b}\,\mathfrak q_{L_\Sigma}[\Gamma_\Sigma^{\mathbf b}]
        -\tfrac14\|\beta_\Sigma\|_{L^\infty}^2\|\Gamma_\Sigma^{\mathbf b}\|_{L^2}^2 ;
    \]
    \item[\textup{(C3)}] \emph{Second-variation saturation (Section~\ref{sec:CylinderRayleigh}):} Hypothesis~\textup{(H6)} (Cauchy--Schwarz saturation~\eqref{eq:CauchySchwarzSaturation}) and Hypothesis~\textup{(H7)} (Jacobi residual control~\eqref{eq:JacobiResidual}) both hold, together with the small-drift regime of Theorem~\ref{thm:PerturbativeDriftStability}.
\end{enumerate}
Moreover, (C2) is automatically satisfied in each of the following sub-regimes:
\begin{enumerate}
    \item[\textup{(C2a)}] \emph{Perturbative:} The collar geometry is $C^{2,\alpha}$-close to a product cylinder with $\|\beta_\Sigma\|_{L^\infty}\le\underline\lambda^{1/2}\varepsilon_{\mathbf b}$ and $\|\mathrm{tf}_\Sigma\,k\|_{W^{1,\infty}}\le\underline\lambda\,\varepsilon_{\mathbf b}$ for $\varepsilon_{\mathbf b}$ sufficiently small (Theorem~\ref{thm:PerturbativeDriftStability});
    \item[\textup{(C2b)}] \emph{Non-perturbative with controlled zero structure:} Conditions (H1)--(H5) of Theorem~\ref{thm:NonPerturbativeDriftStability} hold, i.e.\ the drift-zero locus $Z=\{\beta_\Sigma=0\}$ is a smooth 1-submanifold, $\tr_\Sigma k\in L^\infty$, $\lambda_1^Z(L_\Sigma;\delta)\ge\mu_Z>0$, and the off-$Z$ smallness~\eqref{eq:OffZSmallness} and on-$Z$ bound~\eqref{eq:OnZBound} hold;
    \item[\textup{(C2c)}] \emph{Pure trace:} $\beta_\Sigma\equiv0$ on $\Sigma$ (equivalently $k|_\Sigma$ is pure trace), in which case the original Green-trace calibration of Section~\ref{sec:GreenTraceCalibration} applies.
\end{enumerate}
\end{theorem}

\begin{proof}
(C1) $\Rightarrow$ Penrose is the standard direct MOTS argument of Section~\ref{app:LogicalStructure}. (C2a) $\Rightarrow$ Penrose is Theorem~\ref{thm:PenroseDriftAnchor}. (C2b) $\Rightarrow$ Penrose is Theorem~\ref{thm:NonPerturbativeDriftStability}. (C2c) $\Rightarrow$ Penrose is the original Green-trace calibration (the drift-stability inequality is vacuous when $\beta_\Sigma\equiv0$, and the Green-trace sign reduces to the unanchored case controlled by Theorem~\ref{cor:PenroseFromDualCalibration} under the weighted Riccati positivity conjecture, which in the pure-trace case is automatic because $\tr_\Sigma k = -H_\Sigma$ and the bracket $\mathcal B_\Sigma$ has a definite sign). (C3) $\Rightarrow$ Penrose is Theorem~\ref{thm:GreenSignSecondVariation} combined with Theorem~\ref{thm:PenroseDriftAnchor}.
\end{proof}

\begin{remark}[What is proven unconditionally]\label{rem:UnconditionalResults}
The following are proven unconditionally by the three-channel program:
\begin{enumerate}
    \item The algebraic drift-anchor decomposition of Lemma~\ref{lem:DriftAnchorDecomposition}, which expresses the Green-trace pairing as a drift transport plus a controlled commutator, with the momentum-constraint constant $c_0$ and the residual $\mathfrak G_\Sigma$ explicitly identified.
    \item The nonnegativity of the cylinder Rayleigh quotient $\mathcal R_\Sigma\ge0$ (Lemma~\ref{lem:RayleighNonneg}).
    \item The second-variation positivity identity~\eqref{eq:SecondVariationIdentity}--\eqref{eq:JacobiExpansion} of Theorem~\ref{thm:SecondVariationIdentity}, which identifies the quadratic trace term $\mathfrak Q_\Sigma^{(2)}\ge0$ as a structural source of positivity invisible to first-variation arguments.
    \item The structure of the drift-zero locus $Z=\{\beta_\Sigma=0\}$ as a stratified set of Hausdorff dimension $\le1$ (Lemma~\ref{lem:DriftZeroStructure}).
    \item The Faber--Krahn lower bound~\eqref{eq:FaberKrahnLower} on the localized stability spectrum $\lambda_1^Z(L_\Sigma;\delta)$.
    \item The perturbative drift-stability theorem (Theorem~\ref{thm:PerturbativeDriftStability}) covering collars $C^{2,\alpha}$-close to a product cylinder.
    \item The no-go theorem of Theorem~\ref{thm:NoGoVariational}, which shows that (C1) cannot be derived from the adjoint-cone Euler condition alone.
\end{enumerate}
\end{remark}

\begin{remark}[What is conditional]\label{rem:ConditionalResults}
The following are conditional on hypotheses not yet proven in full generality:
\begin{enumerate}
    \item Hypothesis~\ref{hyp:DriftStability} (drift-stability inequality) outside the perturbative regime (C2a) and the controlled-zero-structure regime (C2b). The remaining cases are those where the drift-zero locus $Z$ has singularities or where the off-$Z$ smallness~\eqref{eq:OffZSmallness} fails.
    \item Hypothesis~\textup{(H6)} (Cauchy--Schwarz saturation~\eqref{eq:CauchySchwarzSaturation}) outside the small-drift regime. This is the condition under which the quadratic trace term $\mathfrak Q_\Sigma^{(2)}$ can be harnessed by the drift-anchor decomposition.
    \item Hypothesis~\textup{(H7)} (Jacobi residual control~\eqref{eq:JacobiResidual}). This controls the non-quadratic-trace part of the Jacobi integral, specifically the $\partial_\nu H + V_\Sigma$ term and the cross-terms $\mathcal J_\Sigma^{\mathrm{cross}}$.
    \item The weighted Riccati positivity conjecture of Section~\ref{sec:GreenTraceCalibration} (Conjecture~\ref{cor:PenroseFromDualCalibration}), which is the original Green-trace route's open condition in the non-pure-trace case.
\end{enumerate}
\end{remark}

\begin{remark}[Logical independence of the three channels]\label{rem:LogicalIndependence}
The three channels (C1), (C2), (C3) are logically independent:
\begin{itemize}
    \item (C1) is a pointwise sign condition on $\tr_\Sigma k$;
    \item (C2) is an integrated drift-weighted inequality;
    \item (C3) is a quadratic saturation condition combining the cylinder Fredholm theory with the second variation of the null expansion.
\end{itemize}
None implies another in general. A complete non-time-symmetric proof would close any one of them in full generality. The drift-anchor decomposition (Lemma~\ref{lem:DriftAnchorDecomposition}) and the second-variation positivity identity (Theorem~\ref{thm:SecondVariationIdentity}) are unconditional structural results that any future proof along this route can build on.
\end{remark}

\subsection{Honest summary}

\begin{remark}[State of the program]\label{rem:StateOfProgram}
The spacetime Penrose inequality along the Jang-equation route is proven in the following regimes: time-symmetric initial data ($k\equiv0$, classical Bray--Khuri); perturbatively near-product collars (Theorem~\ref{thm:PerturbativeDriftStability}); non-perturbatively with controlled zero structure of the drift-zero locus (Theorem~\ref{thm:NonPerturbativeDriftStability}); and the pure-trace case (where the drift vanishes and the original Green-trace calibration applies). The remaining open cases are: general non-time-symmetric initial data where the drift-zero locus has singularities or where the drift transport is not small off $Z$. The three-channel program identifies the precise analytic obstructions in each case and provides unconditional structural identities (the drift-anchor decomposition and the second-variation positivity identity) that any future proof must either use or replicate by other means.
\end{remark}
